% ============================================================================
% COURS 11 — APPLICATIONS GUIDÉES ISSUES DU DOCUMENT D'ORIGINE
% ============================================================================

\section{Applications guidées et exemples technologiques}

\subsection{Pilote automatique de bateau : modèle technologique}

L'ensemble gouvernail--barre franche 1 est en liaison pivot d'axe
$(A,\vec z)$ avec la coque 0. Lorsque le pilotage est automatique,
l'ensemble 1 est actionné par un vérin constitué d'une tige 2 et d'un
corps 3.

\begin{center}
\includegraphics[width=.86\linewidth]{image5.png}
\end{center}

Pour une étude purement cinématique, la liaison entre 1 et 0 peut être
modélisée par une liaison pivot unique. Pour dimensionner les composants,
le modèle doit être plus proche de la réalisation technologique : les
deux roulements sont représentés séparément par une rotule de centre
$A_1$ et une liaison linéaire annulaire de centre $A_2$, de direction
$\vec z$.

\begin{center}
\includegraphics[width=.88\linewidth]{image6.png}
\end{center}

Le modèle d'actions mécaniques doit également faire apparaître :
\begin{itemize}
  \item l'action de l'eau sur le gouvernail, modélisée par une force
        $\vec F_{\mathrm{eau}\rightarrow1}=-F\vec x$ appliquée au point $P$ ;
  \item les actions de pression du fluide sur la tige et le corps du vérin ;
  \item la pesanteur sur les ensembles dont la masse n'est pas négligée.
\end{itemize}

\begin{center}
\includegraphics[width=.90\linewidth]{image7.png}
\end{center}

\begin{TLMethod}{Isolement du vérin complet $\Sigma=2+3$}
Les actions entre la tige 2 et le corps 3 deviennent intérieures. Les
actions extérieures restantes sont les actions de 0 sur 3, de 1 sur 2
et, si elles ne sont pas négligées, les actions de pesanteur sur 2 et 3.
Les torseurs doivent ensuite être transportés au même point avant leur
sommation.
\end{TLMethod}

\subsection{Grue de port : stratégie d'isolement}

On cherche à déterminer l'effort du vérin sur le bras 2 permettant de
maintenir une grue en équilibre sous l'effet de la pesanteur.

\begin{center}
\includegraphics[width=.48\linewidth]{image13.png}
\end{center}

La stratégie efficace est la suivante :
\begin{enumerate}
  \item isoler d'abord la bielle ou le sous-ensemble soumis à deux
        glisseurs, afin de déterminer immédiatement la direction de son
        action ;
  \item isoler ensuite le bras 2, qui est alors soumis à un nombre
        d'inconnues compatible avec les équations du PFS ;
  \item choisir un point de moment qui élimine les réactions de la
        liaison avec le bâti ;
  \item projeter le théorème du moment suivant l'axe normal au plan pour
        obtenir directement l'effort du vérin.
\end{enumerate}

\begin{TLWarning}{Nombre d'inconnues}
Une liaison possédant $I_c$ inconnues cinématiques transmet, dans le
modèle parfait, $I_s=6-I_c$ inconnues statiques. Une liaison hélicoïdale
ne possède qu'une mobilité indépendante : la translation et la rotation
sont liées par le pas.
\end{TLWarning}

\subsection{Poussoir et coulisseau : préparer la résolution}

Le mécanisme comporte un poussoir 1 en rotation par rapport au bâti et
un coulisseau 2 maintenu en contact avec le poussoir. Le poussoir est
soumis à un couple moteur $C_m$ et le coulisseau à l'action d'un fluide.

\begin{center}
\includegraphics[width=.43\linewidth]{image17.jpeg}
\end{center}

On associe au bâti le repère
$R_0=(O,\vec x_0,\vec y_0,\vec z_0)$ et au poussoir le repère
$R_1=(O,\vec x_1,\vec y_1,\vec z_1)$. Le point de contact $A$ est repéré
par $\overrightarrow{BA}=\lambda\vec y_0$.

L'objectif est d'établir une relation entre la force $F$ appliquée au
coulisseau et le couple moteur $C_m$ lorsque le mécanisme est en
équilibre.

\begin{TLMethod}{Stratégie proposée}
\begin{enumerate}
  \item Isoler le coulisseau 2 pour déterminer la direction et
        l'intensité de l'action de contact de 1 sur 2.
  \item Utiliser le théorème des actions réciproques pour obtenir
        l'action de 2 sur 1.
  \item Isoler le poussoir 1.
  \item Appliquer le théorème du moment statique au point $O$ suivant
        $\vec z_0$ : les réactions de la liaison pivot sont éliminées et
        l'équation relie directement $F$ et $C_m$.
\end{enumerate}
\end{TLMethod}

\subsection{Vérin différentiel : action de la pression}

Un vérin possède un piston de diamètre $D_p=\SI{50}{\milli\metre}$ et
une tige de diamètre $D_t=\SI{18}{\milli\metre}$. La chambre arrière est
alimentée sous une pression $p_H$ ; la chambre avant reste alimentée
sous une pression $p_0=\SI{6}{\bar}$.

Les surfaces utiles sont :
\[
S_p=\frac{\pi D_p^2}{4},
\qquad
S_t=\frac{\pi D_t^2}{4}.
\]

L'action globale de la pression sur la tige est un glisseur porté par
l'axe du vérin :
\[
\left\{\mathcal T_{p\rightarrow\mathrm{tige}}\right\}=
\left\{
\begin{array}{c}
\left[p_HS_p-p_0(S_p-S_t)\right]\vec y_3\\
\vec0
\end{array}
\right\}.
\]

Pour un effort utile $F_v=\SI{220}{\newton}$ :
\[
p_H=\frac{F_v+p_0(S_p-S_t)}{S_p}.
\]
L'application numérique conduit à une pression voisine de
$\SI{6.7}{\bar}$.

\subsection{Rendement et puissances perdues}

Pour une transmission recevant une puissance $\mathcal P_e>0$ et
fournissant une puissance utile $\mathcal P_s>0$, le rendement est :
\[
\eta=\frac{\mathcal P_s}{\mathcal P_e},
\qquad 0<\eta\leq1.
\]

La puissance dissipée vaut :
\[
\mathcal P_{\mathrm{perdue}}
=\mathcal P_e-\mathcal P_s
=(1-\eta)\mathcal P_e.
\]

Lorsqu'une transmission est utilisée en sens inverse, il ne faut pas
inverser arbitrairement le rendement. On repart de la définition en
identifiant la puissance entrante et la puissance sortante dans la
situation étudiée.

\subsection{Questions de synthèse}

\begin{enumerate}
  \item Expliquer pourquoi un modèle cinématique minimal peut être
        insuffisant pour dimensionner les roulements d'un mécanisme.
  \item Sur un graphe de structure, montrer comment distinguer les
        actions intérieures et extérieures après un isolement.
  \item Justifier le choix d'un point de réduction qui élimine une action
        inconnue dans une équation de moment.
  \item Établir le torseur global d'une pression uniforme appliquée sur
        une surface plane.
  \item Pour une chaîne comportant plusieurs solides, proposer les
        isolements successifs permettant de déterminer un effort
        d'actionneur avec le minimum d'équations.
\end{enumerate}
